% Section IV -- What Over-Training Reveals (v2). Modules ra2_wedge (+ review), ra1_modelfree (Farseer), ra3_econ
% (aggregate trend), ra5_theory (Props. 2, 4, 5). Numbers and sources: paper/notes/writer2_wedge.md.
\section{What Over-Training Reveals}\label{sec:wedge}

This section uses the technology estimated in Section~\ref{sec:tech} to read the choices of the developers who release
models. A developer that trains a model on more tokens per
parameter than the training-optimal ratio gives up training efficiency to obtain a smaller model of the same quality. By
Proposition~\ref{prop:wedge} it does so only if compactness has value to it, and the technology converts the size of the
sacrifice into that value. Under our reference technology, the median open-weight base model released in 2023--2025 was
built as if the value of compactness were three quarters of lifetime cost. The level rests on an extrapolation that no
public design observes; for most models, the sign does not.

\subsection{From Tokens per Parameter to the Value of Compactness}\label{sec:wedge-measure}

Proposition~\ref{prop:wedge} gives the first-order conditions of a developer that chooses parameters $N$ and training
tokens $D$ under any objective $V(L,N)$ that falls with loss $L$ and may depend directly on size. When the price of a
training FLOP does not vary with model size, data cost nothing beyond compute and no constraint binds, the wedge
$w\equiv\varepsilon_N/\varepsilon_D$ satisfies $w-1=m_N$, the marginal value of compactness at fixed quality per unit of training
expenditure. In the $\kappa$ family the wedge depends on a model's inputs only through its tokens per parameter,
$M=D/N$, and its compute, $C=6ND$ (Online Appendix Corollary~\ref{cor:suff}):
\begin{equation}\label{eq:sIV-inv}
\ln w=\Big(\frac1{\sigma^*}-1\Big)\ln\frac{M}{M^*(C)},\qquad
s\equiv\frac{w-1}{w}=1-\Big(\frac{M}{M^*(C)}\Big)^{-(1/\sigma^*-1)}.
\end{equation}
The data supply $M$ and $C$. The technology supplies two numbers: the zero point $M^*(C)$, the training-optimal ratio at
the model's own compute, and the scale $1/\sigma^*-1$, which turns log distance from the expansion path into a log wedge.
The irreducible loss $E$, the outer exponent $\kappa$ and Hicks-neutral productivity do not enter
(Proposition~\ref{prop:A-wedge}(ii)).

Everything ordinal in equation~\eqref{eq:sIV-inv} is in $M/M^*(C)$. At given compute every technology in the family ranks
models by $M$, and in the sample below a regression of $\ln\hat w$ under the reference technology on $\ln M$ alone has an
$R^2$ of 0.9999. That small models are trained on far more tokens per parameter than large ones can be read from model
cards. What the technology adds is the translation of $M$ into an economic quantity, and this section is organized
around defending the two numbers that make it.

We report the expenditure share $s=m_N/(1+m_N)$, the value of compactness as a share of training expenditure plus that
value, which needs no conversion of FLOPs into dollars. For a developer that serves its own model and bears the cost,
$s$ is the planned serving share of lifetime cost; for an open-weight developer whose users bear the serving cost, it is
an internalized value of compactness, through adoption, memory tiers or latency, in the same units
(Section~\ref{sec:framework:wedge}). A count of lifetime serving tokens $T$ also requires the price ratio $p$ of a
serving to a training FLOP: at the median wedge reported below, $T/D=3(w-1)/p$ is 9.0 at $p=1$ and 0.9 at $p=10$.

Proposition~\ref{prop:modelfree} identifies $w$ without a functional form, but only where a design varies inputs locally
in two directions. Released models lie outside every design, in $M$ and in $C$, so equation~\eqref{eq:sIV-inv} is an
extrapolation. Section~\ref{sec:wedge-pi} checks its direction in $M$ and bounds it in $C$.

\subsection{Sample and Technologies}\label{sec:wedge-data}

\textit{Models.}---The clean inference-demand sample is defined ex ante, before any wedge was computed. It starts from the
173 open-weight base models of the verified sample, whose parameter and token counts are documented by their developers
(Online Appendix~\ref{app:data}). Nine sequential steps then drop models whose $(N,D)$ is not a choice of the kind
Proposition~\ref{prop:wedge} describes: continued-pretraining re-releases (one observation per pretraining run); research
suites and replications whose $D$ was fixed by design, such as Pythia, OPT and BLOOM; mixture-of-experts models; models
that were distilled, pruned or initialized from another model; multimodal pretraining; releases without a base
checkpoint; non-transformer architectures; and token counts without a primary source. The result is 77 models from 36
families and 18 developers, released 2023--2025. Online Appendix~\ref{app:wedge} documents every step and its source;
along the way the median reference wedge stays between 3.77 and 4.46.

\textit{Technologies.}---An ex-ante rule admits 32 technologies: every technology estimated on a designed sweep in
Section~\ref{sec:tech} (each corpus and parameter convention, with $\kappa=1$ and with $\kappa$ free), the path of each
IsoFLOP design combined with its model-free curvature, the literature's estimates on the Chinchilla data, and the
published allocation laws of DeepSeek LLM \citep{bi2024deepseek} and MiniCPM \citep{hu2024minicpm}. The reference is the
Chinchilla technology with $\kappa$ free. Its curvature, $\sigma^*_\kappa=0.701$, matches the model-free estimate of
Section~\ref{sec:tech}, and its zero point is $M^*(10^{21})=22.7$. Its intervals come from a design-conditional wild
bootstrap of the technology \citep{feng2011wild}, and statements about the sample, such as medians and shares, use the
joint distribution across bootstrap draws. Imposing $\kappa=1$, which Section~\ref{sec:tech} rejects, would lower the
median wedge from 3.99 to 3.34.

Where a developer has published its own scaling evidence, its technology is primary for its models. For Meta this is its
IsoFLOP path, which reproduces its published law \citep{grattafiori2024llama}, with the model-free curvature of Meta's own
IsoFLOP profiles ($\sigma^*=0.660$). For AI2 it is the OLMo ladder with $\kappa$ free \citep{bhagia2024establishing}; for
DeepSeek, its published law with the reference curvature; and for Marin, its ladder path with its model-free curvature
($\sigma^*=0.705$; \citealt{marin2026ladders}). These lab-own technologies cover 22 models. For 11 of them the lab's law
predates the model; for the other 11 (Llama~1 and~2, OLMo~1, Marin 8B) it is ex post. Each technology is evaluated in its
own parameter convention, with the output head counted in FLOPs where parameters exclude embeddings. The convention
matters for sub-billion-parameter models, but embeddings are at most 28 percent of parameters in the clean sample, and
the sample median moves little (Online Appendix~\ref{app:wedge}).

\subsection{Revealed Shares}\label{sec:wedge-results}

\begin{figure}[tp]
\centering
\includegraphics[width=\textwidth]{ra2_wedge_ecdf}
\caption{The Distribution of Revealed Expenditure Shares across Technologies}\label{fig:wedge}
\begin{figurenotes}
Empirical distribution functions of $s=(w-1)/w$ across models. Panel a: the 77 models of the clean inference-demand
sample. Panel b: the 22 models covered by a lab-own technology (Meta, AI2, DeepSeek, Marin). Gray: each of the 32 ex-ante
technologies. Blue: the reference technology (Chinchilla, $\kappa$ free). Orange: lab-own technology where available,
otherwise the reference. Values of $s$ below $-1$ ($w<0.5$) are drawn at $-1$. For developers that serve their models,
$s$ is the planned serving share of lifetime cost; for others, the internalized value of compactness in the same units.
\end{figurenotes}
\end{figure}

Figure~\ref{fig:wedge} plots the distribution of $s$ across the clean sample, and Table~\ref{tab:wedge} reports selected
models. Under the reference technology, the median model has $\hat w=3.99$ and $s=0.75$ (95 percent joint bootstrap
interval [0.69, 0.79]), and 97 percent of models are over-trained ($\hat w>1$). The median model was built as if a
model one percent smaller at the same quality were worth three percent of its training cost. For a developer that serves
the model, this is a planned serving share of three quarters of lifetime cost.

The interval measures sampling error within one technology; the choice of technology matters more. Across the 32 ex-ante
technologies the median $s$ ranges from 0.17 to 0.85 (the gray lines in Figure~\ref{fig:wedge}). The spread has a
pattern. Technologies that let the data choose the curvature give medians of 0.69 to 0.79: the $\kappa$-free Chinchilla
and Farseer fits, Meta's and Marin's paths with model-free curvature, and DeepSeek's law with the reference curvature.
The lowest medians come from $\kappa=1$ curvatures far above the model-free estimates, such as those of Marin's
Approach-2 laws ($\sigma^*$ of 0.87 to 0.90), or from a zero point far above the rest, such as MiniCPM's $M^*(10^{21})=192$, which lies
outside MiniCPM's own design. The highest, up to 0.85, come from the $\kappa$-free fits of the small sweeps, whose
curvatures ($\sigma^*_\kappa$ of 0.51 to 0.62) are heterogeneous. Every technology puts the point estimate above one for 64
percent of models, but the union of the 32 technologies' 95 percent intervals lies above one for only 18 percent.

Curvature governs levels, not signs. Holding the reference zero point fixed, lowering $\sigma^*$ from 0.80 to 0.60 raises
the median wedge from 2.25 to 8.64 and the median $s$ from 0.55 to 0.88. The share of models with $\hat w>1$ does not
move, because it depends only on $M^*(C)$. This is why the model-free $\sigma^*$ of about 0.70 from
Section~\ref{sec:tech}, and not a $\kappa=1$ form, should set the scale.

% Table 3 (Section IV): revealed wedges and expenditure shares for selected models.
% Condensed from output/tables/ra2_wedge_models_selected.tex (module ra2_wedge, reviewed run); every cell checked
% against output/tables/ra2_wedge_models.csv (columns N, D, M, M_over_Mstar, w_chin_q, wlo_chin_q, whi_chin_q,
% s_ref, w_lab, s_lab, band_lo, band_hi, serve, wtotalN_chin_q). Sources: paper/notes/writer2_wedge.md.
\begin{table}[tp]
\centering
\caption{Revealed Wedges and Expenditure Shares for Selected Models}
\label{tab:wedge}
\footnotesize
\setlength{\tabcolsep}{1.9pt}
\ifdefined\wedgetabbox\else\newsavebox{\wedgetabbox}\fi
\sbox{\wedgetabbox}{%
\begin{tabular}{@{}lcrrrrr@{\hspace{2pt}}lrcl@{}}
\toprule
 & & $N$ & $D$ & & & \multicolumn{2}{c}{Reference} & & Lab-own & Band \\
Model & Serves & (B) & (T) & $M$ & $M/M^*$ & $\hat w$ & {\footnotesize[95\%]} & $s$ & $\hat w$ ($s$) & of $\hat w$ \\
\midrule
\multicolumn{11}{@{}l}{\textit{Panel A. Clean inference-demand sample}} \\[1pt]
Llama 3 8B & Y & 8.03 & 15.0 & 1,868 & 86.4 & 6.73 & {\footnotesize[5.16, 8.84]} & 0.85 & 8.36 (0.88) & {\footnotesize[0.47, 168]} \\
Llama 3 70B & Y & 70.6 & 15.0 & 213 & 10.0 & 2.67 & {\footnotesize[2.05, 3.48]} & 0.63 & 2.51 (0.60) & {\footnotesize[0.17, 11.7]} \\
Llama 3.1 405B & Y & 406 & 15.6 & 38 & 1.8 & 1.29 & {\footnotesize[0.99, 1.69]} & 0.23 & 0.97 ($-$0.03) & {\footnotesize[0.04, 8.2]} \\
Llama 2 70B & N & 69.0 & 2.0 & 29 & 1.3 & 1.13 & {\footnotesize[0.93, 1.37]} & 0.12 & 0.97$^{e}$ ($-$0.03) & {\footnotesize[0.07, 4.1]} \\
Qwen2.5 7B & Y & 7.62 & 18.0 & 2,364 & 109 & 7.44 & {\footnotesize[5.66, 9.87]} & 0.87 & -- & {\footnotesize[0.49, 231]} \\
Qwen2.5 72B & Y & 72.7 & 18.0 & 248 & 11.7 & 2.86 & {\footnotesize[2.18, 3.74]} & 0.65 & -- & {\footnotesize[0.17, 13.4]} \\
Qwen3 0.6B & Y & 0.60 & 36.0 & 60,398 & 2,759 & 29.6 & {\footnotesize[21.1, 43.0]} & 0.97 & -- & {\footnotesize[1.8, 15,362]} \\
Qwen3 14B & Y & 14.8 & 36.0 & 2,438 & 114 & 7.57 & {\footnotesize[5.60, 10.2]} & 0.87 & -- & {\footnotesize[0.38, 218]} \\
DeepSeek LLM 7B & Y & 7.00 & 2.0 & 286 & 13.0 & 2.99 & {\footnotesize[2.49, 3.60]} & 0.67 & 3.42 (0.71) & {\footnotesize[0.43, 15.8]} \\
DeepSeek LLM 67B & Y & 67.0 & 2.0 & 30 & 1.4 & 1.15 & {\footnotesize[0.94, 1.39]} & 0.13 & 1.34 (0.26) & {\footnotesize[0.07, 4.1]} \\
Gemma 2 27B & Y & 27.2 & 13.0 & 477 & 22.3 & 3.77 & {\footnotesize[2.92, 4.84]} & 0.73 & -- & {\footnotesize[0.27, 29.5]} \\
OLMo 2 7B & N & 7.30 & 4.05 & 555 & 25.4 & 3.99 & {\footnotesize[3.21, 4.94]} & 0.75 & 4.05 (0.75) & {\footnotesize[0.44, 35.5]} \\
OLMo 2 32B & N & 32.2 & 6.15 & 191 & 8.9 & 2.54 & {\footnotesize[2.01, 3.20]} & 0.61 & 1.28 (0.22) & {\footnotesize[0.23, 9.7]} \\
Marin 8B & N & 8.03 & 12.7 & 1,582 & 73.1 & 6.26 & {\footnotesize[4.83, 8.16]} & 0.84 & 5.08$^{e}$ (0.80) & {\footnotesize[0.46, 133]} \\
SmolLM2 1.7B & N & 1.71 & 11.0 & 6,428 & 293 & 11.3 & {\footnotesize[8.69, 15.1]} & 0.91 & -- & {\footnotesize[0.96, 864]} \\
SmolLM2 135M & N & 0.13 & 2.0 & 14,868 & 658 & 16.0 & {\footnotesize[12.4, 21.2]} & 0.94 & -- & {\footnotesize[1.9, 2,701]} \\
\addlinespace
\multicolumn{11}{@{}l}{\textit{Panel B. Mixture-of-experts models, excluded from the sample (bounds at total and active $N$)}} \\[1pt]
DeepSeek-V3 & Y & \multicolumn{4}{l}{685/37 B, 14.8 T} & \multicolumn{2}{l}{[1.01, 3.50]} & \multicolumn{2}{l}{[0.01, 0.71]} & \\
Kimi K2 & Y & \multicolumn{4}{l}{1,026/32 B, 15.5 T} & \multicolumn{2}{l}{[0.87, 3.79]} & \multicolumn{2}{l}{[$-$0.15, 0.74]} & \\
Qwen3 235B-A22B & Y & \multicolumn{4}{l}{235/22 B, 36.0 T} & \multicolumn{2}{l}{[2.34, 6.39]} & \multicolumn{2}{l}{[0.57, 0.84]} & \\
Qwen3 30B-A3B & Y & \multicolumn{4}{l}{30.5/3.3 B, 36.0 T} & \multicolumn{2}{l}{[5.56, 14.3]} & \multicolumn{2}{l}{[0.82, 0.93]} & \\
\bottomrule
\end{tabular}%
}
\ifdim\wd\wedgetabbox>\textwidth\resizebox{\textwidth}{!}{\usebox{\wedgetabbox}}\else\usebox{\wedgetabbox}\fi\par
\medskip
\begin{tablenotes}
Serves: Y if the developer served its own language models through a first-party API or product when the model was
released. $N$: total parameters; $D$: documented training tokens; $M=D/N$; $M^*$: the reference technology's
training-optimal ratio at the model's own compute. Reference: Chinchilla with $\kappa$ free ($\sigma^*_\kappa=0.701$);
$\hat w$ from equation~(\ref{eq:sIV-inv}); brackets: 95 percent design-conditional wild-bootstrap interval of one
technology, not the uncertainty of the wedge. $s=(w-1)/w$, negative when $w<1$. Lab-own: Meta, Meta's IsoFLOP path with
its model-free curvature; DeepSeek, its published law; AI2, the OLMo ladder with $\kappa$ free; Marin, its ladder path
with its model-free curvature. $^{e}$Lab law postdates the model. Band: union of the 95 percent intervals (points for
published laws) over the 32 ex-ante technologies.
\end{tablenotes}
\begin{tablenotes}[Source]
Per-model values under every technology and convention are in the replication package; see Online
Appendix~\ref{app:wedge}.
\end{tablenotes}
\end{table}


Lab-own technologies give similar answers. On the 22 models they cover, the median $s$ is 0.61 under the lab's own
technology and 0.65 under the reference (Figure~\ref{fig:wedge}, panel b). Meta is the sharpest case. Under Meta's own
law and curvature, Llama~3 8B has $\hat w=8.4$ (95 percent interval [5.7, 12.8]) and $s=0.88$. Meta served Llama~3
through its own assistant, so this is a planned serving share. The 405-billion-parameter flagship of the same herd has
$\hat w=0.97$; that it lies on Meta's training-optimal path is by construction, since Meta sized it with this law.
Within one developer and one data pipeline, then, the flagship was trained as if for training efficiency and the
smallest model as if for serving. The level for the small model depends on the technology Meta is taken to have used:
across Meta's own alternatives, the 8B's wedge runs from 3.1 under the $\kappa=1$ law to 12.2 with the path refitted on
the budgets whose IsoFLOP minima are bracketed.

\subsection{Extrapolation and Partial Identification}\label{sec:wedge-pi}

Released models lie beyond the designs in two directions.

\textit{Tokens per parameter.}---Forty-seven of the 77 models have $M$ above 341, the largest ratio in the Chinchilla
design. Inside that range the median $s$ is 0.56, and no model has all 32 technologies' intervals above one; beyond it the
median is 0.85. Only two models lie inside the design in both $M$ and compute. The agreement among technologies at
high $M$ therefore owes much to a shared functional form, under which $\hat w\to\infty$ as $M\to\infty$ at given
compute.

Farseer's near-factorial design \citep{li2025predictableb}, which reaches $M=2{,}570$ in non-embedding parameters,
allows a direct check. Proposition~\ref{prop:modelfree}(i) identifies the wedge at each grid point from local slopes of
the loss surface, without a functional form. Figure~\ref{fig:farseer} compares this local wedge with those of parametric
forms fitted to the same runs. Every form understates the wedge at high $M$, and the gap grows with $M$. At
$M\ge1{,}024$, where the local wedge is 7.5, the $\kappa$-free form's wedge is 0.70 to 0.78 times the local one (the range
over kernel bandwidths) and the Chinchilla form's is 0.44 to 0.50 times. The Chinchilla form fitted only on runs with
$M\le100$, the extrapolation that equation~\eqref{eq:sIV-inv} performs for small over-trained models, gives 0.26 to 0.30
times. Panel c shows the source: $\ln w$ is convex in $\ln(M/M^*(C))$, whereas the $\kappa$ family makes it linear
(Proposition~\ref{prop:modelfree}(iii)), and linearity is rejected. Applying Farseer's gaps to the 57 clean models inside
Farseer's range of $M$ raises their median $s$ under the $\kappa$-free Farseer technology from 0.68 to 0.71 or 0.72,
depending on the bandwidth.

\begin{figure}[tp]
\centering
\includegraphics[width=\textwidth]{ra1_modelfree_farseer_w}
\caption{Parametric Extrapolation against the Model-Free Wedge}\label{fig:farseer}
\begin{figurenotes}
Farseer's 404 released runs \citep{li2025predictableb}. The local wedge is
$w=(L_c-L_x)/(L_c+L_x)$ (Proposition~\ref{prop:modelfree}(i)), from kernel-weighted local quadratic fits of loss on
$(\ln N,\ln D)$, evaluated at grid points inside each model size's token range. Panels a and b: mean of
$\ln(w_{\mathrm{param}}/w_{\mathrm{local}})$ by bin of $M$ for parametric forms fitted to the same runs (Chinchilla form
on all runs and on runs with $M\le100$; $\kappa$ family on all runs), with $N$ excluding (a) or including (b) embeddings;
bars are 95 percent intervals from a wild cluster bootstrap by model size (Webb weights, 999 draws) and exclude
smoothing bias. Panel c: $\ln w$ at grid points against $\ln M$, locally and under the Chinchilla form and Farseer's own
non-homothetic form.
\end{figurenotes}
\end{figure}

Where it can be checked, then, parametric extrapolation in $M$ errs on the conservative side: it understates the wedge
relative to the true technology. Three limits apply. The evidence comes from one recipe, and at the highest ratios from
the smallest models in the design, of 0.10 to 0.34 billion parameters. Nothing is observed beyond Farseer's range, where
20 clean models lie. And revealed preference recovers the technology a developer believed when it chose $(N,D)$
(Proposition~\ref{prop:A-wedge}(ii)): if developers planned with fitted parametric laws, the parametric wedge is the one
their first-order conditions equate, and a correction toward the true technology would move $\hat w$ away from it
(Remark~\ref{rem:A-pi-param}). \textcolor{red}{[TBD-m9: Our controlled experiment repeats the check at ratios up to
$M\approx$~X in non-embedding parameters (pre-registered question Q3): fitted on $M\le100$, the Chinchilla form
over/understates the local wedge at $M>1{,}000$ by a factor of Y (Section~\ref{sec:tech}).]}

\textit{Compute.}---The zero point is the harder problem. No design observes $M^*(C)$ above $10^{22}$ FLOP, while 88
percent of the clean sample used more compute than the largest Chinchilla run ($1.3\times10^{22}$ FLOP).
Proposition~\ref{prop:pi} replaces a point extrapolation with an identified set. We anchor $M^*$ at the largest budget
of each IsoFLOP design: 22.7 at $3\times10^{21}$ FLOP for Chinchilla, 22.3 at $10^{22}$ for Llama~3, 9.4 to 12.4 at $3\times10^{20}$ for Marin's three
corpora, and DeepSeek's published law. We then let the elasticity of $M^*$ with respect to compute, $e=1-2a$, lie
anywhere in the range spanned by these paths and the published laws, $[-0.16,0.19]$. The identified set for
$M^*(10^{24})$ is then [2.3, 89]. With Meta's own anchor alone it narrows to [10.1, 57].

What the data identify at frontier scale is a sign, not a level. For any technology with single-peaked IsoFLOP profiles,
which hold whenever $0<\sigma<1$ (Lemma~\ref{lemma:sigma}), the wedge exceeds one exactly when $M$ exceeds $M^*(C)$
(Proposition~\ref{prop:pi}(iii)). For 86 percent of the clean sample, $M$ exceeds every compute-optimal ratio in the
identified set, so over-training is identified with no assumption on curvature or functional form. With each lab's own
anchor the share is 77 percent. Admitting every ex-ante technology as an anchor, with the widest range of path slopes, it
falls to 23 percent. No model is identified as under-trained under any of these assumptions. The 11 models whose sign is
not identified under the IsoFLOP anchors all have $M\le96$, and they include most large flagships: Llama~2 70B,
Llama~3.1 405B, Qwen2 72B, Falcon 180B and DeepSeek LLM 67B. Magnitudes need a curvature. With $1/\sigma^*-1$ between
0.40 and 0.52, the range of the $\kappa$-free and model-free estimates, the medians of the models' lower and upper bounds
on $s$ are 0.59 and 0.93.

\subsection{Conduct: Who Values Compactness}\label{sec:wedge-conduct}

Proposition~\ref{prop:wedge} admits several objectives, and they predict different patterns
(Remark~\ref{rem:A-conduct}). If developers value compactness because they bear serving costs, the wedge should rise with
a developer's own serving footprint, and providers of closed models, which bear their serving cost in full, should
over-train at least as much as open-weight developers. If open developers value compactness through adoption by users who
bear the cost, the wedge should rise with the sensitivity of adoption to users' cost. If sizes are set by a menu of memory
tiers, sizes should bunch at tier boundaries and the wedge should be unrelated to demand given the tier. Inference below
uses a restricted wild cluster bootstrap by developer \citep{cameron2008bootstrap,webb2023reworking}.

Open-weight models are more over-trained than closed ones. Among 186 models of production scale from 67 developers,
the open-weight premium in $\ln\hat w$ is 0.52 given release year and compute ($p=0.011$). A developer that bears all of its
serving cost should not over-train less than one that bears part of it, so serving costs alone do not explain open
developers' choices. The comparison is limited to disclosed allocations, and closed developers have disclosed few since
2023.

Within the clean sample the tests have little power. The coefficient on the developer serving its own models is 0.45
(standard error 0.26), the sign the serving story predicts, but with 18 developer clusters, 7 of them serving, the
bootstrap $p$-value is 0.39. Stated on-device targets are carried by two developers and cannot be tested. Sizes do bunch:
51 percent of the 220 open-weight production-scale models released in 2023--2026 sit in five windows around quantized
memory tiers, against 13 percent under a smooth counterfactual \citep{chetty2011adjustment,kleven2016bunching}. Given
compute, however, a model's position relative to the tiers predicts its wedge no better than placebo tier menus do
(placebo $p$-values between 0.37 and 1.00). Tiers shape the choice of $N$; the tier model's prediction for the wedge is
not visible.

Because the data do not discriminate between these objectives, we report the estimand separately by who serves. The 33
clean models whose developers served their own language models at release, from 7 developers including Meta, Alibaba,
Google and DeepSeek, have a median $s$ of 0.80, and every one is over-trained. For them $s$ reads as planned serving
expenditure. The other 44 models, from 12 developers including AI2, Hugging Face and TII, have a median $s$ of 0.74. For
them $s$ is an internalized value of compactness, not serving expenditure.

Common token budgets need separate treatment. Thirty-two clean models belong to 11 families that train every size on
one token budget, so their member wedges are not revealed preferences (Proposition~\ref{prop:wedge}(iv)). The
family-level shares that these families reveal run from 0.18 for DeepSeek LLM to 0.89 for Qwen3 (Online Appendix
Proposition~\ref{prop:A-family}). Excluding the 11 families leaves the median $s$ at 0.75, and families with
common and size-specific budgets have similar median wedges (3.90 and 4.01).

The cost side matters little for $s$. If the price of a FLOP, in training and serving alike, varies with model
size with an elasticity between $-0.1$ and $0.1$, the median $s$ is 0.72 to 0.77. If serving cost scales as $N^\phi$ with
$\phi$ between 0.5 and 1.25, it is 0.70 to 0.86, and it never depends on $p$. The rival wedges of
Section~\ref{sec:framework:wedge} have known signs. The two that concern markup estimation, factor-biased recipes and a
mismatch between the loss we measure and the output a developer values
\citep{demirer2020production,raval2023testing,bond2020unpleasant}, are reduced by lab-own technologies for the 22 models
they cover, and the spread across the 32 technologies indicates their size for the rest. Dropping the nine models trained
mainly on synthetic data, whose generation is compute planned at training time, leaves the median $s$ at 0.74.

\subsection{Validation and Trends}\label{sec:wedge-trends}

Levels cannot be validated with public data. Per-model token volumes on OpenRouter, the natural test, are not available,
and one year of developer-level traffic \citep{aubakirova2026state} equals 0.04 to 4.6 percent of the lifetime tokens
that the same developers' releases imply at $p=1$, which says only that the platform carries a small part of a model's
lifetime use (Online Appendix~\ref{app:wedge-trends}). Hugging Face derivative counts give a weaker check. At given compute, models with higher $M$ have more
quantizations, fine-tunes and other derivatives (elasticity 0.80, $p=0.031$; 74 models from 16 developers). Smaller models
of equal compute are adapted more, as planned demand predicts but also as accessibility predicts, so this does not
validate levels.

The trend is clearer than the level. Among open-weight releases of production scale (141 models with at least
$10^{21}$ FLOP; Table~\ref{tab:econ}, panel~C), the compute-weighted share $\sum_i(w_i^+-1)C_i/\sum_iw_i^+C_i$, with $w_i^+=\max(w_i,1)$, rose from 0.04
for releases in 2019--2022 to 0.35 in 2023, 0.54 in 2024 and 0.82 in 2025. It is planned serving expenditure for
developers that serve and an internalized value of compactness for the rest. The early value reflects two under-trained
models, OPT-175B and GLM-130B, with 92 percent of that period's compute. On the clean sample the aggregate share rises
from 2023 to 2025 under every ex-ante technology; the level does not travel, and for 2025 it ranges from 0.24 to 0.96
across technologies. Llama~3.1 405B carries 39 percent of
2024 compute, and without it the 2024 share is 0.63. Part of the rise may reflect more inference tokens per query, from
reasoning and repeated sampling \citep{roberts2026test}, rather than more queries; both are serving demand.
Section~\ref{sec:econ} compares these shares with firms' disclosures.

Two findings of this section are robust: over-training, identified at frontier scale under weak assumptions for 86
percent of the clean sample, and the steep rise between 2022 and 2025 in the share of lifetime cost that developers
attached to compactness. The level, about three quarters for the median model, is conditional on a zero point that no
public design observes at the relevant compute.
